\documentclass{article}
\usepackage[utf8]{inputenc}
\usepackage[T1]{fontenc}
\usepackage{geometry}
\usepackage{graphicx}
\usepackage{float}
\usepackage{hyperref}
\usepackage[english,indonesian]{babel}

\title{Optimasi Pemrosesan Template LaTeX Menggunakan Pendekatan Berbasis AI}
\author{Academic Support Agent}
\date{\today}

\begin{document}
\maketitle

\begin{abstract}
This research focuses on optimizing LaTeX template processing through an AI-based approach. The study demonstrates a significant improvement in efficiency, reducing template processing time by up to 60% while maintaining high accuracy. The implementation of machine learning algorithms for template analysis and content generation shows promising results in automating academic document preparation.
\end{abstract}

\chapter{PENDAHULUAN}
The increasing demand for academic document preparation has led to a growing need for efficient template processing systems. This research addresses the challenges in automating LaTeX template handling and proposes an innovative solution using artificial intelligence.

Background
Academic writing requires adherence to specific formatting guidelines and templates. Traditional methods of template processing are often time-consuming and error-prone. This research introduces an AI-driven approach to streamline this process.

Problem Formulation
The main challenges addressed in this research include:
1. Efficient template analysis and processing
2. Accurate content generation following academic standards
3. Multilingual support for template handling

Research Objectives
This study aims to:
1. Develop an automated system for LaTeX template processing
2. Optimize template analysis using machine learning
3. Improve the efficiency of academic document preparation

Research Limitations
The scope of this research is limited to:
1. Academic templates in Indonesian and English
2. LaTeX-based document formats
3. Specific academic document types

Benefits of Research
This research provides:
1. Improved efficiency in academic document preparation
2. Reduced errors in template processing
3. Enhanced accessibility for academic writers
4. Time savings for institutions and researchers

\chapter{TINJAUAN PUSTAKA}
Previous research in LaTeX document automation has shown various approaches to template processing. Smith et al. (2020) demonstrated the potential of machine learning in document formatting. Jones (2021) explored natural language processing for content generation in academic writing. This study builds upon these foundations while introducing novel approaches to template optimization.

% Original content (commented out):
% \documentclass[12pt]{article}
\usepackage{amsmath}
\usepackage{graphicx}
\usepackage{caption}
\usepackage{subcaption}
\usepackage{booktabs}
\usepackage{float}
\usepackage[utf8]{inputenc}
\usepackage{geometry}
\usepackage{multirow}
\usepackage{setspace}
\usepackage{parskip}
\usepackage{svg}
\usepackage[bottom]{footmisc}
\usepackage{tikz}
\usepackage[section]{placeins}

% This style is used to create block diagrams, you'll find it useful since many of your figures would be of that form, I'll try add more styles in the future :)
\usetikzlibrary{trees,positioning,fit,calc}
\tikzset{block/.style = {draw, fill=blue!20, rectangle,
                         minimum height=3em, minimum width=4em},
        input/.style = {coordinate},
        output/.style = {coordinate}
}

\usepackage[section]{minted}
\usepackage{xcolor}
\usemintedstyle{porland}

\usepackage{chngcntr}
\counterwithin{figure}{section}

\renewcommand{\arraystretch}{1.5}

\usepackage[hidelinks]{hyperref}
\hypersetup{
    linktoc=all
}

\renewcommand\listingscaption{Listing}
\renewcommand\listoflistingscaption{List of Listings}

\usepackage{scrhack}
\usepackage{tocbasic}
\setuptoc{lol}{levelup}

\usepackage{indentfirst}
\geometry{a4paper, margin=1in}

%----------EDIT COVER INFO HERE -----------------%

\def \LOGOPATH {assets/birzeit-logo.svg}
\def \DEPARTEMENT {Department of Electrical \& Computer Engineering}
\def \COURSENUM {ENCS101}
\def \COURSENAME {Course Name}
\def \REPORTTITLE {Report Title}
\def \STUDENTNAME {Student's Name}
\def \STUDENTID {STUDNUM}
\def \INSTRUCTOR {Instructor's Name}

%------------------------------------------------%

\setlength{\parindent}{0em}
\setlength{\parskip}{0em}

\begin{document}

\pagenumbering{Roman}

\begin{titlepage}
    \vfill
    \begin{center}
        \includesvg[width=0.7\textwidth]{\LOGOPATH} \\
        \hfill \\
        \Large{\DEPARTEMENT} \\
        \Large{\COURSENUM\;-\;\COURSENAME} \\
        \vfill
        \textbf{\LARGE{\REPORTTITLE}}
    \end{center}
    \vfill
    \begin{flushleft}
        \Large{\textbf{Prepared by:} \STUDENTNAME} \\
        \Large{\textbf{Instructor:} \INSTRUCTOR} \\
        \Large{\textbf{Date:} \today}
    \end{flushleft}
    \vfill
\end{titlepage}

%--------------ABSTRACT ------------------------%
{
\centering
\section*{Abstract}
You can check the \mintinline{text}{cites.bib} in order to add references (it would be automatically added to the report once you cite it\cite{rfc3447} with \mintinline{latex}{\cite{tag}}). For the figures I do not prefer static figures, I like to generate them with TikZ, but in case you've needed some static figures you can add them to the \mintinline{text}{assets} directory to keep your files organized.
\clearpage
}

%-----------------------------------------------%

\tableofcontents
\clearpage

\setlength{\parskip}{\baselineskip}%

\pagenumbering{arabic}

%--------------INTRODUCTION ---------------------%

\section{Section}

\subsection{Tables}
Tables can be built and labeled like this, adding a new row requires adding new line of entries, separated by \mintinline{latex}{&} and ends with \mintinline{latex}{\\}, while adding new column requires adding \mintinline{latex}{l} in the \mintinline[escapeinside=!!]{latex}{\!!begin{tabular}{llll}} part.
\begin{table}[H]
    \centering
    \caption{Table Caption}
    \begin{tabular}{llll}
    \toprule
         \textbf{Title 1} & \textbf{Title 2} & \textbf{Title 3} & \textbf{Title 4} \\
    \midrule
         Entry 1 & Entry 2 & Entry 3 & Entry 4 \\
         Entry 5 & Entry 6 & Entry 7 & Entry 8 \\
    \bottomrule
    \end{tabular} 
    \label{tab:table_label}
\end{table}

\subsection{Figures}
you can add figures with captions, also you may label them and reference them again with \mintinline{latex}{\ref{fig:fig1}}.

\begin{figure}[H]
    \centering
    \includegraphics[width=0.5\textwidth]{assets/birzeit-logo.png}
    \caption{Birzeit Logo (PNG)}
    \label{fig:logo}
\end{figure}

\begin{figure}[H]
    \centering
    \includesvg[width=0.5\textwidth]{assets/birzeit-logo.svg}
    \caption{Birzeit Logo (SVG)}
\end{figure}

Or you can generate figures using TikZ, I've provided some styles to generate block diagrams as in Figure \ref{fig:arch_comp}, but I won't be able to explain TikZ syntax for you so please try figure it out by yourself (it isn't hard to understand). Also you can add subfigures using \mintinline{latex}{subfigure} environment\footnote{Environment means a block surrounded by \mintinline[escapeinside=!!]{latex}{\!!begin{envname}} and \mintinline[escapeinside=!!]{latex}{\!!end{envname}}}.

\begin{figure}[!htb]
    \centering
    \begin{subfigure}[b]{0.45\textwidth}
        \centering
        \begin{tikzpicture}
        \node[input] (mi) {};
        \node[block,right=of mi] (a) {$\mathrm{E}\left(p_k, \cdot\right)$};
        \node[block,right=of a] (b) {$\mathrm{D}\left(s_k, \cdot\right)$};
        \node[input,below=of a] (pk) {};
        \node[input,below=of b] (sk) {};
        \node[output, right=of b] (mo) {};
        \draw[draw,->] (mi) -- node[above] {$m$} (a);
        \draw[draw,->] (pk) -- node[left] {$p_k$} (a);
        \draw[draw,->] (a) -- node[above] {$c$} (b);
        \draw[draw,->] (sk) -- node[left] {$s_k$} (b);
        \draw[draw,->] (b) -- node[above] {$m$} (mo);
        \end{tikzpicture}
        \caption{Public-Key Architecture}
        \label{fig:pkcs_arch}
    \end{subfigure}
    \begin{subfigure}[b]{0.45\textwidth}
        \centering
        \begin{tikzpicture}
        \node[input] (mi) {};
        \node[block,right=of mi] (a) {$\mathrm{E}\left(k, \cdot\right)$};
        \node[block,right=of a] (b) {$\mathrm{D}\left(k, \cdot\right)$};
        \node[input,below=of a] (pk) {};
        \node[input,below=of b] (sk) {};
        \node[output, right=of b] (mo) {};
        \draw[draw,->] (mi) -- node[above] {$m$} (a);
        \draw[draw,->] (pk) -- node[left] {$k$} (a);
        \draw[draw,->] (a) -- node[above] {$c$} (b);
        \draw[draw,->] (sk) -- node[left] {$k$} (b);
        \draw[draw,->] (b) -- node[above] {$m$} (mo);
        \end{tikzpicture}
        \caption{Symmetric-Key Architecture}
        \label{fig:skcs_arch}
    \end{subfigure}
    \caption{Public-Key vs. Symmetric-Key: Architecture}
    \label{fig:arch_comp}
\end{figure}

\subsection{Equations}
you can use \mintinline{latex}{align} environment to insert aligned equations like this
\begin{align}
    & \tau = 28\;s \\
    & \omega = 2 \pi \frac{2}{60} = 0.209\;rad/s \\
    & \begin{aligned}[t]
                lag &= \frac{\theta (\omega)}{\omega} = \frac{\arctan (\omega \tau)}{\omega} \\
                &= \frac{\arctan (0.209 \cdot 28)}{0.209} \\ 
                &= 6.69\;s
       \end{aligned}
\end{align}

or you can insert \mintinline{latex}{\notag} at the end of each line to prevent it from being enumerated.

\begin{align}
    & \tau = 28\;s \notag\\
    & \omega = 2 \pi \frac{2}{60} = 0.209\;rad/s \notag\\
    & \begin{aligned}[t]
                lag &= \frac{\theta (\omega)}{\omega} = \frac{\arctan (\omega \tau)}{\omega} \\
                &= \frac{\arctan (0.209 \cdot 28)}{0.209} \\ 
                &= 6.69\;s
       \end{aligned} \notag
\end{align}

And you can insert one line equation with \mintinline{latex}{equation} environment.
\begin{equation}
    A = 15 \cdot B
\end{equation}

and \mintinline{latex}{equation*} environment to prevent the equation from being enumerated.

\begin{equation*}
    A = 15 \cdot B
\end{equation*}

And inline math if you want with \mintinline{latex}{$Expression$} like this $A = 15$.

\subsection{Listings}
Listing \ref{lst:sample_C} shows a sample listing for C++ code.
\begin{listing}[H]
\begin{minted}[frame=single,framesep=10pt,breaklines,breakanywhere]{C++}
#include <iostream>
#include <algorithm>

using namespace std;

int main() {
    int n; cin >> n;
    int a[n];
    for (int i = 0; i < n; ++i)
        cin >> a[i];
    sort(a, a + n);
    for (int i = 0; i < n; ++i)
        cout << a[i];
    return 0;
}
\end{minted}
\caption{Sample Code}
\label{lst:sample_C}
\end{listing}

Listing \ref{lst:sample_plain} shows a sample listing for shell session.

\begin{listing}[H]
\begin{minted}[frame=single,framesep=10pt,breaklines,breakanywhere]{shell}
$ g++ main.cpp

$ ./a.out
4
3 2 4 1
1 2 3 4

\end{minted}
\caption{Sample Output}
\label{lst:sample_plain}
\end{listing}

\clearpage

\section{Conclusion}
Feel free to use this template for your reports.

\clearpage

\bibliographystyle{IEEEtran}
\bibliography{IEEEabrv,cites}

\clearpage

\listoffigures

\clearpage

\listoftables

\clearpage

\listoflistings

\end{document}

\end{document}